\documentclass[11pt,compress,t,notes=noshow, xcolor=table]{beamer}

\input{../../style/preamble}
\input{../../latex-math/basic-math}
\input{../../latex-math/basic-ml}
\input{../../latex-math/optim-basics}

% Lipschitz constant of the Hessian (Boyd's L; plain L is already the smoothness
% constant in this lecture, see 01-foundations/03-convex.tex)
\newcommand{\LH}{L_{\H}}

\title{Optimization in Machine Learning}

\begin{document}

\titlemeta{
Second order methods
}{
Deep Dive: Convergence of Newton-Raphson
}{
figure/NR_phases.png
}{
\item Assumptions
\item Damped phase
\item Quadratically convergent phase
\item Complexity bound
\vspace{0.5cm}
\item[] \hspace{-0.7cm}Slides based on\\
\hspace{-0.7cm}\furtherreading{BOYD2004} (Ch. 9.5.3)
}

\begin{framei}{Setting}
\item \textbf{Recall:}
\begin{itemizeM}
\item with full steps ($\alpha^{[t]} = 1$), NR may overshoot/diverge far from
$\xvs$ $\Rightarrow$ damping is necessary
\item many objective functions are only locally quadratic
\end{itemizeM}
\item In this deep dive, we show that damped NR has two phases:
\begin{itemizeM}
\item \textbf{Damped phase:} quadratic model unreliable, $\alpha^{[t]} < 1$ possible; $f$ decreases by a
\textit{constant} $\delta$ per step
\item \textbf{Pure Newton phase:} quadratic model accurate, $\alpha^{[t]} = 1$ always accepted; \textit{quadratic} convergence
\end{itemizeM}
\splitV[0.62]{
\imageC[0.85]{figure/NR_phases.png}
}{
{\footnotesize \textbf{Example 2} (cont'd):
$$f(\xv) = (1 + x_1^2)^{1/2} + \frac{x_2^2}{2}$$
from $\xv^{[0]} = (3, 2)^\top$,\\
with $\gamma_1 = 0.1$, $\tau = 0.7$}
}
\end{framei}


\begin{framei}[fs=small]{Setting}
\item \textbf{Assumptions} on the sublevel set $\S = \{\xv : f(\xv) \le f(\xv^{[0]})\}$:
\begin{itemizeM}
\item $f \in \CC{2}$, $m$-strongly convex ($\H(\xv) \succeq m \id$) and $L$-smooth ($\H(\xv) \preceq L \id$)
\item $\H$ is \textbf{Lipschitz} with constant $\LH$ ($\rightarrow$ bound on the \textit{3rd} derivative):
$$\|\H(\xv) - \H(\xtil)\|_2 \le \LH \|\xv - \xtil\| \qquad \text{for all } \xv, \xtil \in \S$$
Intuition: $\LH$ measures how well the quadratic model approximates $f$ ($\LH = 0$ for quadratic $f$)
\end{itemizeM}
\item Damped NR with backtracking: $\alpha_{\text{init}} = 1$, $\gamma_1 \in (0, 1/2)$, $\tau \in (0, 1)$
\begin{framed}
\textbf{Theorem \textnormal{(Convergence of damped NR)}.}
With the scaled gradient norm $\theta^{[t]} = \frac{\LH}{2m^2}\|\nabla f(\xv^{[t]})\|$, there exist
$\eta \in (0, m^2/\LH]$ and $\delta > 0$ such that
\begin{align*}
\|\nabla f(\xv^{[t]})\| \ge \eta \; &\Rightarrow \; f(\xv^{[t+1]}) - f(\xv^{[t]}) \le - \delta
&& \textit{(damped)} \\
\|\nabla f(\xv^{[t]})\| < \eta \; &\Rightarrow \; \alpha^{[t]} = 1 \, \text{ and } \,
\theta^{[t+1]} \le \big( \theta^{[t]} \big)^2 && \textit{(pure)}
\end{align*}
\end{framed}
\end{framei}

% \begin{framei}{Outline}
% \item We start by analyzing the implications of the theorem (main convergence result)
% \item We then actually prove the theorem
% \begin{itemizeM}
% \item \textbf{Damped phase:} We show that each step reduces $f$ by at least a constant $\delta$ and show how $\delta$ looks like
% \item \textbf{Pure phase:} We show that $\alpha^{[t]} = 1$ is always accepted and that the scaled gradient norm $\theta^{[t]}$ decreases quadratically
% \end{itemizeM}
% \end{framei}

\begin{framei}[fs=small]{What the theorem implies: pure phase}
\item [] Before we actually prove the theorem, we analyze its implications for the convergence of damped NR.
\item \textbf{Message 1:} once NR enters the pure phase, it stays there
\begin{itemizeM}
  \item suppose we are in the pure phase at step $t$, i.e. $\|\nabla f(\xv^{[t]})\| < \eta$
  \item since $\eta \le m^2/\LH$ by def., the scaled gradient norm is below $\tfrac{1}{2}$:
  $$\theta^{[t]} = \frac{\LH}{2m^2}\|\nabla f(\xv^{[t]})\|
  < \frac{\LH}{2m^2} \cdot \frac{m^2}{\LH} = \frac{1}{2}$$
  \item with this and the theorem:
  $$\theta^{[t+1]} \le \big( \theta^{[t]} \big)^2 = \theta^{[t]} \cdot \theta^{[t]}
  < \tfrac{1}{2} \theta^{[t]}$$
  $$\Rightarrow \|\nabla f(\xv^{[t+1]})\| < \tfrac{1}{2} \|\nabla f(\xv^{[t]})\| < \eta$$
  \item i.e., the condition is also satisfied at step $t+1$ and by recursion at all future steps $l \geq t$, implying (by the theorem):
  $$\alpha^{[l]} = 1 \quad \text{ and } \quad \theta^{[l+1]} \le \big( \theta^{[l]} \big)^2$$
\end{itemizeM}
\end{framei}

\begin{framei}[fs=small]{What the theorem implies: pure phase}
\item \textbf{Message 2:} convergence is extremely rapid in the pure phase
\begin{itemizeM}
\item $\theta^{[t+1]} \le \big( \theta^{[t]} \big)^2$ is exactly \textbf{quadratic convergence}
\item Applying it recursively, for $l \geq t$:
$$\theta^{[l]} \le \big( \theta^{[t]} \big)^{2^{l-t}} \leq \left(\frac{1}{2}\right)^{2^{l-t}}$$
\item Strong convexity gives $f(\xv) - f(\xvs) \le \frac{1}{2m}\|\nabla f(\xv)\|^2$ and thus:
$$f(\xv^{[l]}) - f(\xvs) \le \eps_0 \big( \theta^{[l]} \big)^2
\le \eps_0 \left(\frac{1}{2}\right)^{2^{l-t+1}}, \qquad \eps_0 = \frac{2m^3}{\LH^2}$$
\item Inverting this for the step count yields:\footnote{\footnotesize Exactly
$\log_2 \log_2 (\eps_0/\eps) - 1$ steps; the $-1$ is dropped for readability.}
\begin{framed}
$$f(\xv^{[l]}) - f(\xvs) \le \eps \quad \text{after at most} \quad \log_2 \log_2 (\eps_0/\eps)
\; \text{ pure steps}$$
\end{framed}
\item[$\Rightarrow$] \textbf{Note}: at most 6 pure steps give an accuracy of
$\eps \approx 5 \cdot 10^{-20} \, \eps_0$
\end{itemizeM}
\end{framei}

\begin{framei}[fs=small]{What the theorem implies: damped phase \& final result}
\item \textbf{Message 3:} the number of steps in the damped phase is bounded
\begin{itemizeM}
\item Every damped step reduces $f$ by at least $\delta$, so after $T$ damped steps
$f(\xv^{[T]}) \le f(\xv^{[0]}) - T \delta$
\item Since $f$ cannot decrease below $f(\xvs)$:
\begin{framed}
$$\#\{\text{damped steps}\} \le \frac{f(\xv^{[0]}) - f(\xvs)}{\delta}$$
\end{framed}
\item[$\Rightarrow$] finite, but expensive phase; grows with the initial gap
\end{itemizeM}
\item \textbf{Final result} (both phases together): to reach $f(\xv^{[t]}) - f(\xvs) \le \eps$:
\begin{framed}
$$\#\text{iterations} \;\; \le \;\;
\underbrace{\frac{f(\xv^{[0]}) - f(\xvs)}{\delta}}_{\text{damped}} \; + \;
\underbrace{\log_2 \log_2 (\eps_0/\eps)}_{\text{pure}, \; \approx \, 6}$$
\end{framed}
\end{framei}


\begin{framev}[fs=small]{Proof strategy}
\textbf{Note:} we have now seen the implications of the theorem, but have not yet proven it, i.p. the existence of $\eta$ and $\delta$. Our proof strategy is as follows:
\begin{enumerate}
\item \textbf{Tool:} the \textbf{Newton decrement} $\lambda(\xv)$ as progress measure
\item \textbf{Damped phase:} $L$-smoothness bounds the step size from below
\begin{framed}
$$\alpha^{[t]} \ge \tau \frac{m}{L}
\quad \Longrightarrow \quad
\delta = \gamma_1 \tau \eta^2 \frac{m}{L^2}$$
\end{framed}
\item \textbf{Pure phase:} Lipschitz $\H$ keeps the quadratic model accurate enough that
\begin{framed}
$$\alpha^{[t]} = 1 \text{ is accepted}
\quad \text{and} \quad
\theta^{[t+1]} \le \big( \theta^{[t]} \big)^2$$
$$\text{for} \quad \eta = \min\{1, \, 3(1 - 2\gamma_1)\} \, \frac{m^2}{\LH}$$
\end{framed}
\end{enumerate}
\end{framev}


\begin{framei}[fs=small]{Tool: the Newton decrement}
\item Backtracking needs $\nabla f(\xv)^\top \mathbf{d}$ along $\mathbf{d} = - \H(\xv)^{-1} \nabla f(\xv)$;
its negative defines the \textbf{Newton decrement} $\lambda(\xv)$:
\begin{framed}
$$\lambda(\xv)^2 = \underbrace{\nabla f(\xv)^\top \H(\xv)^{-1} \nabla f(\xv)}_{\text{definition}}
= \underbrace{\mathbf{d}^\top \H(\xv) \mathbf{d}}_{\text{squared } \H\text{-norm of } \mathbf{d}}
= \underbrace{- \nabla f(\xv)^\top \mathbf{d}}_{-\, \text{dir. derivative}}$$
\end{framed}
\item[$\Rightarrow$] the Armijo rule for the Newton step reads
$f(\xv + \alpha \mathbf{d}) \le f(\xv) - \gamma_1 \alpha \lambda(\xv)^2$
\item[$\Rightarrow$] $\lambda$ measures the gradient in the geometry of $\H$, not the Euclidean one
\item From $m \id \preceq \H(\xv) \preceq L \id$, i.e. $\tfrac{1}{L}\id \preceq \H(\xv)^{-1} \preceq \tfrac{1}{m}\id$:
\begin{framed}
$$\frac{1}{L}\|\nabla f(\xv)\|^2 \le \lambda(\xv)^2 \le \frac{1}{m}\|\nabla f(\xv)\|^2,
\qquad \|\mathbf{d}\| \le \frac{\lambda(\xv)}{\sqrt{m}}$$
\end{framed}
\end{framei}


\begin{framei}[fs=small]{Damped phase: the step size cannot get small}
\item Assume $\|\nabla f(\xv^{[t]})\| \ge \eta$ and abbreviate $\mathbf{d} = \mathbf{d}^{[t]}$,
$\lambda = \lambda(\xv^{[t]})$
\item \textbf{Recall} the descent lemma ($L$-smoothness) along $\mathbf{d}$, for any $\alpha \ge 0$:
$$f(\xv^{[t]} + \alpha \mathbf{d}) \le f(\xv^{[t]}) + \alpha \nabla f(\xv^{[t]})^\top \mathbf{d}
+ \frac{L}{2}\alpha^2 \|\mathbf{d}\|^2
\le f(\xv^{[t]}) - \alpha \lambda^2 + \frac{L}{2m}\alpha^2 \lambda^2$$
using $\nabla f(\xv^{[t]})^\top \mathbf{d} = - \lambda^2$ and $\|\mathbf{d}\|^2 \le \lambda^2/m$
\item $\hat\alpha = m/L$ minimizes the right side, and because $\gamma_1 < \tfrac{1}{2}$:
$$f(\xv^{[t]} + \hat\alpha \mathbf{d}) \le f(\xv^{[t]}) - \frac{m}{2L}\lambda^2
\le f(\xv^{[t]}) - \gamma_1 \hat\alpha \lambda^2$$
\item[$\Rightarrow$] $\hat\alpha$ passes the Armijo rule, and so does every smaller step; backtracking
therefore accepts before shrinking past $\tau \hat\alpha$:
\begin{framed}
$$\alpha^{[t]} \ge \tau \frac{m}{L}$$
\end{framed}
\end{framei}


\begin{framei}[fs=small]{Damped phase: guaranteed decrease}
\item The accepted step satisfies the Armijo rule, hence
\begin{align*}
f(\xv^{[t+1]}) - f(\xv^{[t]}) &\le - \gamma_1 \alpha^{[t]} \lambda^2 \\
&\le - \gamma_1 \tau \frac{m}{L} \lambda^2 \\
&\le - \gamma_1 \tau \frac{m}{L^2} \|\nabla f(\xv^{[t]})\|^2 \\
&\le - \gamma_1 \tau \frac{m}{L^2} \eta^2
\end{align*}
\item \textbf{Arguments:} Armijo rule (line~1). Step size bound (line~1 $\to$ 2).
$\lambda^2 \ge \|\nabla f\|^2/L$ (line~2 $\to$ 3). Phase condition (line~3 $\to$ 4).
\begin{framed}
$$\delta = \gamma_1 \tau \eta^2 \frac{m}{L^2}$$
\end{framed}
\item[$\Rightarrow$] the decrease does not shrink as $t$ grows, making Message 3 possible
\end{framei}


\begin{framei}[fs=small]{Pure phase: a cubic bound along the Newton step}
\item Now assume $\|\nabla f(\xv^{[t]})\| < \eta$ and restrict $f$ to the Newton ray:
$\tilde{f}(\alpha) = f(\xv^{[t]} + \alpha \mathbf{d})$, so $\tilde{f}'(0) = -\lambda^2$ and
$\tilde{f}''(\alpha) = \mathbf{d}^\top \H(\xv^{[t]} + \alpha \mathbf{d}) \mathbf{d}$
\item By the Lipschitz condition between $\xv^{[t]} + \alpha \mathbf{d}$ and $\xv^{[t]}$, then
sandwiched with $\mathbf{d}$ (Cauchy-Schwarz):
\begin{align*}
\big\| \H(\xv^{[t]} + \alpha \mathbf{d}) - \H(\xv^{[t]}) \big\|_2 &\le \alpha \LH \|\mathbf{d}\| \\
\big| \mathbf{d}^\top \big( \H(\xv^{[t]} + \alpha \mathbf{d}) - \H(\xv^{[t]}) \big) \mathbf{d} \big|
&\le \alpha \LH \|\mathbf{d}\|^3
\end{align*}
\item This is $|\tilde{f}''(\alpha) - \tilde{f}''(0)|$; with $\tilde{f}''(0) = \lambda^2$, $\|\mathbf{d}\| \le \lambda/\sqrt{m}$:
$$\tilde{f}''(\alpha) \le \lambda^2 + \alpha \frac{\LH}{m^{3/2}} \lambda^3$$
\item Integrating twice (with $\tilde{f}'(0) = -\lambda^2$) and setting $\alpha = 1$:
$$f(\xv^{[t]} + \mathbf{d}) \le f(\xv^{[t]}) - \frac{1}{2}\lambda^2 + \frac{\LH}{6 m^{3/2}} \lambda^3$$
\end{framei}


\begin{framei}[fs=small]{Pure phase: the full step is accepted}
\item Strong convexity turns the phase condition into a bound on $\lambda$: if
$\eta \le 3(1 - 2\gamma_1)\frac{m^2}{\LH}$, then
$$\lambda \le \frac{\|\nabla f(\xv^{[t]})\|}{\sqrt{m}} < \frac{\eta}{\sqrt{m}}
\le 3(1 - 2\gamma_1)\frac{m^{3/2}}{\LH}
\quad \Longrightarrow \quad
\frac{\LH \lambda}{6 m^{3/2}} \le \frac{1 - 2\gamma_1}{2}$$
\item Factoring the cubic bound, the bracket is therefore at least $\gamma_1$:
\begin{align*}
f(\xv^{[t]} + \mathbf{d}) &\le f(\xv^{[t]}) - \lambda^2 \underbrace{\left( \frac{1}{2}
- \frac{\LH \lambda}{6 m^{3/2}} \right)}_{\ge \, \gamma_1}
\le f(\xv^{[t]}) - \gamma_1 \lambda^2 \\
&= f(\xv^{[t]}) + \gamma_1 \cdot 1 \cdot \nabla f(\xv^{[t]})^\top \mathbf{d}
\end{align*}
\item[$\Rightarrow$] this is exactly the Armijo rule at $\alpha = 1$: \textbf{backtracking accepts the full step
straight away}
\item[$\Rightarrow$] here $\gamma_1 < \tfrac{1}{2}$ is essential: otherwise $3(1 - 2\gamma_1) \le 0$ and no
such $\eta$ exists
\end{framei}


\begin{framei}[fs=small]{Pure phase: quadratic convergence}
\item The Newton step solves $\nabla f(\xv^{[t]}) + \H(\xv^{[t]}) \mathbf{d} = \zero$, so we can subtract
it:
\begin{align*}
\|\nabla f(\xv^{[t+1]})\| &= \|\nabla f(\xv^{[t]} + \mathbf{d}) - \nabla f(\xv^{[t]}) - \H(\xv^{[t]})\mathbf{d}\| \\
&= \left\| \int_0^1 \left( \H(\xv^{[t]} + \alpha \mathbf{d}) - \H(\xv^{[t]}) \right) \mathbf{d} \; d\alpha \right\|
\; \le \; \frac{\LH}{2}\|\mathbf{d}\|^2 \\
&= \frac{\LH}{2}\|\H(\xv^{[t]})^{-1} \nabla f(\xv^{[t]})\|^2
\; \le \; \frac{\LH}{2m^2}\|\nabla f(\xv^{[t]})\|^2
\end{align*}
\item \textbf{Arguments:} fundamental theorem of calculus (line~1 $\to$ 2). Lipschitz $\H$ and
$\int_0^1 \alpha \, d\alpha = \tfrac{1}{2}$ (line~2). $\|\H^{-1}\| \le 1/m$ (line~3).
\item Multiplying by $\frac{\LH}{2m^2}$ gives the theorem's form; taking the smaller of the two ceilings
on $\eta$ (previous slide and theorem/Message 1):
\begin{framed}
$$\theta^{[t+1]} \le \big( \theta^{[t]} \big)^2, \qquad
\eta = \min\{1, \, 3(1 - 2\gamma_1)\} \, \frac{m^2}{\LH}$$
\end{framed}
\end{framei}


\begin{framei}[fs=small]{Putting everything together}
\item Substituting $\eta$ into $\delta = \gamma_1 \tau \eta^2 \frac{m}{L^2}$:
$$\delta = \gamma_1 \tau \min\{1, \, 9(1 - 2\gamma_1)^2\} \, \frac{m^5}{L^2 \LH^2}$$
\item Inserting this into the final result and using $\log_2 \log_2 (\eps_0/\eps) \approx 6$:
\begin{framed}
$$\#\text{iterations} \; \le \; 6 \; + \;
\frac{L^2 \LH^2 / m^5}{\gamma_1 \tau \min\{1, \, 9(1 - 2\gamma_1)^2\}}
\left( f(\xv^{[0]}) - f(\xvs) \right)$$
\end{framed}
\item the accuracy $\eps$ has collapsed into the constant $6$
\item the cost is driven by how far the
starting point is from the optimum
\item all the cost sits in the damped phase, with
$\frac{L^2 \LH^2}{m^5} = \kappa^2 \big( \frac{\LH}{m^{3/2}} \big)^2$: ill-conditioning ($\kappa = L/m$) and non-quadraticity
\item $\gamma_1, \tau$ enter only mildly $\Rightarrow$ backtracking parameters barely matter
\end{framei}


\begin{framei}[fs=small]{Discussion}
\item Slideset 1 claimed NR converges \textit{locally} quadratically, and may diverge for bad starting
points -- we can now make both precise
\item \textbf{``Locally''} is the pure phase, and the neighborhood is quantified:
$$\|\nabla f(\xv^{[t]})\| < \eta = \min\{1, \, 3(1 - 2\gamma_1)\} \, \frac{m^2}{\LH}$$
\item \textbf{Damping upgrades local to global:} the damped phase provably reaches this neighborhood
after finitely many steps
\item[$\Rightarrow$] damped NR converges from \textit{any} $\xv^{[0]}$ whose sublevel set satisfies the
assumptions -- a stronger statement than slideset 1 made
\item \textbf{Undamped NR} ($\alpha^{[t]} = 1$ always): the pure-phase proof never used backtracking, only
that a \textit{full} step was taken $\Rightarrow$ quadratic convergence still holds \textit{inside} the region
\item[] but the damped-phase argument collapses: nothing forces $f$ to decrease (cf. Example 2)
\end{framei}


\begin{framei}[fs=small]{Discussion: the role of the assumptions}
\item $m$ \textbf{(strong convexity)}: translates between $\|\nabla f\|$, $\lambda$, $\|\mathbf{d}\|$ and
turns a small gradient into a small optimality gap
\item $L$ \textbf{(smoothness)}: descent lemma $\Rightarrow$ backtracking cannot return arbitrarily tiny steps
\item $\LH$ \textbf{(Lipschitz $\H$)}: controls the entire pure phase; for quadratic $f$ we have $\LH = 0$,
so $\eta \to \infty$ and $\theta^{[t]} = 0$ $\Rightarrow$ NR converges in a single step
\item All assumptions are only needed on $\S$, which suffices because $f$ decreases every step
\item \textbf{Caveat:} $m, L, \LH$ are unknown in practice, and none of them is affine invariant --
although NR itself is
\item[$\Rightarrow$] motivates \textit{self-concordance} \furtherreading{BOYD2004} (Ch. 9.6): a convergence
analysis free of unknown constants
\end{framei}

\endlecture
\end{document}
